\newpage
\section{Математическое описание}
\label{sec:math}

\subsection{Конечный автомат}
\label{sec:automaton}

Устройство управления электронными часами описывается конечным автоматом Мили
\[
	M = (Q, X, Y, \delta, \lambda, S_0),
\]
где
\begin{itemize}
	\item $Q = \{S_0,S_1,\ldots,S_{11}\}$~--- множество состояний
	(режимов управления);
	\item $X$ и $Y$~--- множества входных и выходных воздействий;
	\item $\delta\colon Q \times X \to Q$~--- функция переходов;
	\item $\lambda\colon Q \times X \to Y$~--- функция выходов;
	\item $S_0$~--- начальное состояние.
\end{itemize}

В такте $t$ при текущем состоянии $q(t)$ и входном символе $x(t)$ автомат формирует выход $y(t)$ и переходит в следующее состояние:
\[
	y(t) = \lambda(q(t), x(t)), \qquad
	q(t+1) = \delta(q(t), x(t)), \qquad q(0) = S_0.
\]

Состояние автомата задает режим индикации и назначение кнопок. Текущее время, показание и признак хода секундомера, время будильника, признак его включения и номер мелодии хранятся в соответствующих блоках часов.

Граф автомата приведен на рис.~\ref{fig:clock-control}. Метка $x/z_i$ задает вход и выход перехода. Если указан только вход, выход равен $z_0$. Если для кнопки дуга отсутствует, автомат сохраняет состояние и выдает $z_0$.

\begin{figure}[H]
	\centering
	\tikzset{
		state/.style={draw,circle,minimum size=8mm,inner sep=0pt,font=\small},
		event/.style={font=\footnotesize,inner sep=2pt},
		edge/.style={-{Stealth[length=2mm]}}
	}

	\small Корректировка и останов текущего времени\par
	\begin{tikzpicture}
		\node[state] (s0) at (0,0) {$S_0$};
		\node[state] (s1) at (3.4,0) {$S_1$};
		\node[state] (s2) at (6.8,0) {$S_2$};
		\node[state] (s3) at (10.2,0) {$S_3$};
		\node[state] (s10) at (13.6,0) {$S_{10}$};
		\draw[edge] (-1,0) -- (s0);
		\draw[edge] (s0) -- node[event,above] {$a/z_0$} (s1);
		\draw[edge] (s1) -- node[event,above] {$a,c\,/\,z_0$} (s2);
		\draw[edge] (s2) -- node[event,above] {$a,c\,/\,z_0$} (s3);
		\draw[edge] (s3) -- node[event,above] {$a,c\,/\,z_0$} (s10);
		\draw[edge] (s1) to[loop above] node[event,above] {$b/z_1$} (s1);
		\draw[edge] (s2) to[loop above] node[event,above] {$b/z_2$} (s2);
		\draw[edge] (s3) to[loop above] node[event,above] {$b/z_3$} (s3);
		\draw[edge,rounded corners=3pt]
			(s10.south) -- (13.6,-0.95) --
			node[event,below] {$a,b,c\,/\,z_0$} (0,-0.95) -- (s0.south);
	\end{tikzpicture}

	\medskip
	Секундомер\par
	\begin{tikzpicture}
		\node[state] (s0) at (6.6,0) {$S_0$};
		\node[state] (s4) at (6.6,-1.8) {$S_4$};
		\node[state] (s5) at (1.2,-4.05) {$S_5$};
		\node[state] (s6) at (12,-4.05) {$S_6$};
		\draw[edge] (s0) to[bend left=32] node[event,right] {$b/z_0$} (s4);
		\draw[edge] (s4) to[bend left=32] node[event,left] {$c/z_0$} (s0);
		\draw[edge] (s4) to[bend left=10] node[event,below] {$a/z_4$} (s5);
		\draw[edge] (s5) to[bend left=10] node[event,above=2pt] {$a/z_0$} (s4);
		\draw[edge] (s4) to[bend left=10] node[event,above=2pt] {$b/z_5$} (s6);
		\draw[edge] (s6) to[bend left=10] node[event,below] {$b/z_6$} (s4);
		\draw[edge] (s5) to[bend left=5] node[event,above] {$b/z_5$} (s6);
		\draw[edge] (s6) to[bend left=5] node[event,below] {$a/z_4$} (s5);
		\draw[edge] (s5) to[bend left=12] node[event,above left] {$c/z_0$} (s0);
		\draw[edge] (s6) to[bend right=12] node[event,above right] {$c/z_0$} (s0);
	\end{tikzpicture}

	\medskip
	Будильник\par
	\begin{tikzpicture}
		\node[state] (s0) at (0,0) {$S_0$};
		\node[state] (s7) at (3.4,0) {$S_7$};
		\node[state] (s8) at (6.8,0) {$S_8$};
		\node[state] (s9) at (10.2,0) {$S_9$};
		\node[state] (s11) at (13.6,0) {$S_{11}$};
		\draw[edge] (s0) to[bend left=18] node[event,above] {$c/z_0$} (s7);
		\draw[edge] (s7) to[bend left=18] node[event,below] {$c/z_0$} (s0);
		\draw[edge] (s7) -- node[event,above] {$b/z_0$} (s8);
		\draw[edge] (s8) -- node[event,above] {$b,c\,/\,z_0$} (s9);
		\draw[edge] (s9) -- node[event,above] {$b,c\,/\,z_0$} (s11);
		\draw[edge] (s7) to[loop above] node[event,above] {$a/z_7$} (s7);
		\draw[edge] (s8) to[loop above] node[event,above] {$a/z_8$} (s8);
		\draw[edge] (s9) to[loop above] node[event,above] {$a/z_9$} (s9);
		\draw[edge] (s11) to[loop above] node[event,above] {$a/z_{10}$} (s11);
		\draw[edge,rounded corners=3pt]
			(s11.south) -- (13.6,-0.85) --
			node[event,below] {$b,c\,/\,z_0$} (3.4,-0.85) -- (s7.south);
	\end{tikzpicture}

	\medskip
	Четвертьчасовая сигнализация\par
	\begin{tikzpicture}
		\node[state] (s0) at (0,0) {$S_0$};
		\node[state] (s12) at (6.8,0) {$S_{12}$};
		\draw[edge,dashed] (s0) to[bend left=12] node[event,above] {$\tau/z_0$} (s12);
		\draw[edge,dashed] (s12) to[bend left=12] node[event,below] {$\rho/z_0$} (s0);
	\end{tikzpicture}

	\caption{Граф конечного автомата управления}
	\label{fig:clock-control}
\end{figure}


\newpage
\subsubsection{Описание состояний конечного автомата управления}
\label{sec:states}

В каждом состоянии определены показания индикатора и действия кнопок.

\begin{enumerate}[itemsep=0.45\baselineskip]
	\item $S_0$~--- отображение текущего времени. Это начальное состояние. На индикаторе видны часы, минуты и секунды в 24-часовом формате.
	\begin{itemize}
		\item $a$~--- перейти к корректировке секунд ($S_1$).
		\item $b$~--- показать секундомер ($S_4$).
		\item $c$~--- показать время будильника и признак его включения ($S_7$).
		\item $d$~--- сохранить текущее состояние.
	\end{itemize}

	\item $S_1$~--- корректировка секунд. На индикаторе отображается поле секунд.
	\begin{itemize}
		\item $b$~--- увеличить секунды на единицу по модулю $60$, оставаясь в $S_1$.
		\item $a$, $c$, $d$~--- перейти к корректировке минут ($S_2$).
	\end{itemize}

	\item $S_2$~--- корректировка минут. На индикаторе отображается поле минут.
	\begin{itemize}
		\item $b$~--- увеличить минуты на единицу по модулю $60$, оставаясь в $S_2$.
		\item $a$, $c$, $d$~--- перейти к корректировке часов ($S_3$).
	\end{itemize}

	\item $S_3$~--- корректировка часов. На индикаторе отображается поле часов.
	\begin{itemize}
		\item $b$~--- увеличить часы на единицу по модулю $24$, оставаясь в $S_3$.
		\item $a$, $c$, $d$~--- завершить корректировку и остановить ход часов ($S_{10}$).
	\end{itemize}

	\item $S_4$~--- отображение секундомера. На индикаторе видны минуты и секунды отсчета. Признак хода секундомера сохраняется при входе в этот режим из $S_0$.
	\begin{itemize}
		\item $a$~--- запустить или продолжить отсчет с сохраненного значения ($S_5$).
		\item $b$~--- остановить секундомер и зафиксировать показание ($S_6$).
		\item $c$, $d$~--- вернуться к текущему времени ($S_0$), сохранив признак хода секундомера.
	\end{itemize}

	\item $S_5$~--- запущенный секундомер. На индикаторе обновляется показание секундомера. Любая кнопка возвращает автомат в $S_4$ при продолжающемся отсчете.

	\item $S_6$~--- остановленный секундомер. На индикаторе остается зафиксированное показание.
	\begin{itemize}
		\item $c$~--- сбросить показание до $00{:}00$, оставаясь в $S_6$.
		\item $a$, $b$, $d$~--- вернуться в $S_4$ при остановленном отсчете.
	\end{itemize}

	\item $S_7$~--- настройка будильника. На индикаторе видны установленное время будильника и признак его включения.
	\begin{itemize}
		\item $a$~--- переключить признак включения будильника, оставаясь в $S_7$.
		\item $b$~--- перейти к корректировке минут будильника ($S_8$).
		\item $c$~--- показать номер мелодии будильника ($S_{11}$).
		\item $d$~--- вернуться к отображению текущего времени ($S_0$).
	\end{itemize}

	\item $S_8$~--- корректировка минут будильника. На индикаторе отображаются минуты установленного времени.
	\begin{itemize}
		\item $a$~--- увеличить минуты на единицу по модулю $60$, оставаясь в $S_8$.
		\item $b$, $c$, $d$~--- перейти к корректировке часов будильника ($S_9$).
	\end{itemize}

	\item $S_9$~--- корректировка часов будильника. На индикаторе отображаются часы установленного времени.
	\begin{itemize}
		\item $a$~--- увеличить часы на единицу по модулю $24$, оставаясь в $S_9$.
		\item $b$, $c$, $d$~--- вернуться к настройкам будильника ($S_7$).
	\end{itemize}

	\item $S_{10}$~--- останов часов. Индикатор показывает время на момент останова. Нажатие любой кнопки возобновляет ход часов и возвращает автомат в $S_0$.

	\item $S_{11}$~--- выбор мелодии будильника. На индикаторе отображается номер выбранной мелодии.
	\begin{itemize}
		\item $a$~--- выбрать следующую мелодию, оставаясь в $S_{11}$. После последней мелодии выбирается первая.
		\item $b$, $c$, $d$~--- вернуться к настройкам будильника ($S_7$).
	\end{itemize}
\end{enumerate}

Вершина $S_{12}$ исходного рисунка соответствует активному состоянию звукового блока. В начале очередной четверти часа автомат выдает команду $z_{12}$, и блок включает сигнал на одну секунду. По истечении секунды блок возвращается в режим ожидания. Состояние автомата управления и назначение кнопок при этом сохраняются.

Коды состояний автомата управления приведены в таблице~\ref{tab:states}. Для 12 состояний необходимо $\lceil\log_2 12\rceil = 4$ разряда. Код $q_3q_2q_1q_0$ совпадает с двоичной записью номера состояния, $q_3$~--- старший разряд.

\begin{table}[H]
	\centering
	\caption{Состояния и их двоичные коды}
	\label{tab:states}
	\small\setstretch{1.12}
	\renewcommand{\arraystretch}{1.12}
	\begin{tabularx}{\textwidth}{|c|>{\raggedright\arraybackslash}p{6.2cm}|>{\raggedright\arraybackslash}X|c|}
		\hline
		Состояние & Режим & Показание индикатора & Код \\
		\hline
		$S_0$ & Отображение времени & Часы, минуты и секунды & $0000$ \\
		$S_1$ & Корректировка секунд & Секунды & $0001$ \\
		$S_2$ & Корректировка минут & Минуты & $0010$ \\
		$S_3$ & Корректировка часов & Часы & $0011$ \\
		$S_4$ & Отображение секундомера & Минуты и секунды отсчета & $0100$ \\
		$S_5$ & Запущенный секундомер & Текущее показание & $0101$ \\
		$S_6$ & Остановленный секундомер & Зафиксированное показание & $0110$ \\
		$S_7$ & Настройка будильника & Время и признак включения & $0111$ \\
		$S_8$ & Установка минут будильника & Минуты будильника & $1000$ \\
		$S_9$ & Установка часов будильника & Часы будильника & $1001$ \\
		$S_{10}$ & Остановленные часы & Время на момент останова & $1010$ \\
		$S_{11}$ & Выбор мелодии будильника & Номер мелодии & $1011$ \\
		\hline
	\end{tabularx}
\end{table}


\subsubsection{Описание и кодирование входных воздействий}
\label{sec:inputs}

Входы $X = \{a,b,c,\tau,\rho,\varnothing\}$ кодируются тремя разрядами $x_2x_1x_0$ (таблица~\ref{tab:inputs}).

\begin{table}[H]
	\centering
	\caption{Кодирование входных воздействий}
	\label{tab:inputs}
	\small\setstretch{1.12}
	\renewcommand{\arraystretch}{1.12}
	\begin{tabularx}{\textwidth}{|c|>{\raggedright\arraybackslash}X|c|}
		\hline
		Вход & Воздействие & Код \\
		\hline
		$a$ & Нажатие кнопки $a$ & $000$ \\
		$b$ & Нажатие кнопки $b$ & $001$ \\
		$c$ & Нажатие кнопки $c$ & $010$ \\
		$\tau$ & Начало четверти часа (минуты $00$, $15$, $30$, $45$) & $100$ \\
		$\rho$ & Истечение одной секунды после входа в $S_{12}$ & $101$ \\
		$\varnothing$ & Отсутствие события & $110$ \\
		\hline
	\end{tabularx}
\end{table}

Событие $\tau$ формируется при естественном ходе часов в начале очередной четверти часа при нулевом значении секунд. Событие $\rho$ вырабатывает таймер, запущенный при входе в $S_{12}$.

\newpage
\subsubsection{Описание и кодирование выходных воздействий}
\label{sec:outputs}

Выходной алфавит $Y = \{z_0,z_1,\ldots,z_{10}\}$ содержит команды, которые автомат подает при обработке входных событий. Выход $z_0$ означает работу по режиму, заданному состоянием. Например, при переходе в $S_0$ индикатор показывает текущее время, а в $S_{10}$ отображается время на момент останова.

Выходы кодируются четырьмя разрядами $y_3y_2y_1y_0$ по номеру символа (таблица~\ref{tab:outputs}).

\begin{table}[H]
	\centering
	\caption{Кодирование выходных воздействий}
	\label{tab:outputs}
	\small\setstretch{1.12}
	\renewcommand{\arraystretch}{1.10}
	\begin{tabularx}{\textwidth}{|c|>{\raggedright\arraybackslash}X|c|}
		\hline
		Выход & Действие & Код \\
		\hline
		$z_0$ & Работа по режиму, заданному состоянием & $0000$ \\
		$z_1$ & Сбросить значение секунд текущего времени до $00$ & $0001$ \\
		$z_2$ & Увеличить значение минут на единицу & $0010$ \\
		$z_3$ & Увеличить значение часов на единицу & $0011$ \\
		$z_4$ & Разрешить отсчет секундомера с сохраненного значения & $0100$ \\
		$z_5$ & Остановить секундомер и сохранить показание & $0101$ \\
		$z_6$ & Сбросить показание остановленного секундомера & $0110$ \\
		$z_7$ & Переключить признак включения будильника & $0111$ \\
		$z_8$ & Увеличить значение минут будильника на единицу & $1000$ \\
		$z_9$ & Увеличить значение часов будильника на единицу & $1001$ \\
		$z_{10}$ & Выбрать следующую мелодию будильника & $1010$ \\
		\hline
	\end{tabularx}
\end{table}

Команда $z_4$ устанавливает признак хода секундомера, а $z_5$ снимает его. Команда $z_6$ обнуляет показание остановленного секундомера. В $S_{10}$ ход счетчика текущего времени приостановлен. Четвертьчасовой звук разрешен в течение пребывания автомата в $S_{12}$, поэтому обе дуги сигнализации имеют выход $z_0$.

\newpage
\subsubsection{Функции переходов и выходов}
\label{sec:functions}

В таблице~\ref{tab:functions} указаны пары $\delta(q,x)/\lambda(q,x)$ для нажатий кнопок. Строка задает текущее состояние $q$, столбец~--- входной символ $x$.

\begin{table}[H]
	\centering
	\caption{Функции переходов и выходов при нажатии кнопок}
	\label{tab:functions}
	\small\setstretch{1.12}
	\renewcommand{\arraystretch}{1.1}
	\setlength{\tabcolsep}{9pt}
	\begin{tabular}{|c|cccc|}
		\hline
		$q \backslash x$ & $a$ & $b$ & $c$ & $d$ \\
		\hline
		$S_{0}$ & $S_{1}/z_{0}$ & $S_{4}/z_{0}$ & $S_{7}/z_{0}$ & $S_{0}/z_{0}$ \\
		$S_{1}$ & $S_{2}/z_{0}$ & $S_{1}/z_{1}$ & $S_{2}/z_{0}$ & $S_{2}/z_{0}$ \\
		$S_{2}$ & $S_{3}/z_{0}$ & $S_{2}/z_{2}$ & $S_{3}/z_{0}$ & $S_{3}/z_{0}$ \\
		$S_{3}$ & $S_{10}/z_{0}$ & $S_{3}/z_{3}$ & $S_{10}/z_{0}$ & $S_{10}/z_{0}$ \\
		$S_{4}$ & $S_{5}/z_{4}$ & $S_{6}/z_{5}$ & $S_{0}/z_{0}$ & $S_{0}/z_{0}$ \\
		$S_{5}$ & $S_{4}/z_{4}$ & $S_{4}/z_{4}$ & $S_{4}/z_{4}$ & $S_{4}/z_{4}$ \\
		$S_{6}$ & $S_{4}/z_{5}$ & $S_{4}/z_{5}$ & $S_{6}/z_{6}$ & $S_{4}/z_{5}$ \\
		$S_{7}$ & $S_{7}/z_{7}$ & $S_{8}/z_{0}$ & $S_{11}/z_{10}$ & $S_{0}/z_{0}$ \\
		$S_{8}$ & $S_{8}/z_{8}$ & $S_{9}/z_{0}$ & $S_{9}/z_{0}$ & $S_{9}/z_{0}$ \\
		$S_{9}$ & $S_{9}/z_{9}$ & $S_{7}/z_{0}$ & $S_{7}/z_{0}$ & $S_{7}/z_{0}$ \\
		$S_{10}$ & $S_{0}/z_{0}$ & $S_{0}/z_{0}$ & $S_{0}/z_{0}$ & $S_{0}/z_{0}$ \\
		$S_{11}$ & $S_{11}/z_{11}$ & $S_{7}/z_{0}$ & $S_{7}/z_{0}$ & $S_{7}/z_{0}$ \\
		\hline
	\end{tabular}
\end{table}

Для событий времени и такта без события правила одинаковы во всех состояниях $q\in Q$:
\[
	\begin{array}{c|ccc}
		x & \tau & \alpha & \varnothing \\ \hline
		\delta(q,x) & q & q & q \\
		\lambda(q,x) & z_{12} & z_{13} & z_0
	\end{array}
\]

Каждой паре из $Q\times X$ соответствует ровно одно следующее состояние и один выход. Все состояния достижимы из $S_0$ по нажатию кнопок, а из каждого режима можно вернуться в $S_0$.

Для просмотра текущего времени при работающем секундомере из $S_0$ последовательно нажимаются $b$, $a$, $b$, $c$. Автомат проходит состояния $S_4$, $S_5$, $S_4$ и $S_0$. На последнем переходе выдается $z_0$, поэтому секундомер продолжает отсчет. Чтобы снова увидеть показание секундомера и остановить его, нужно дважды нажать $b$.

\newpage
В таблице~\ref{tab:binary-functions} те же переходы по кнопкам записаны в двоичном виде. В каждой ячейке сначала приведен код следующего состояния, а затем код выхода. Для входов $100$, $101$ и $110$ код состояния сохраняется, а код выхода равен соответственно $1100$, $1101$ и $0000$.

\begin{table}[H]
	\centering
	\caption{Двоичные коды переходов и выходов при нажатии кнопок}
	\label{tab:binary-functions}
	\small\setstretch{1.12}
	\renewcommand{\arraystretch}{1.1}
	\begin{tabularx}{\textwidth}{|c|*{4}{>{\centering\arraybackslash}X}|}
		\hline
		Состояние & $000$ & $001$ & $010$ & $011$ \\
		\hline
		$S_{0}$ & $0001/0000$ & $0100/0000$ & $0111/0000$ & $0000/0000$ \\
		$S_{1}$ & $0010/0000$ & $0001/0001$ & $0010/0000$ & $0010/0000$ \\
		$S_{2}$ & $0011/0000$ & $0010/0010$ & $0011/0000$ & $0011/0000$ \\
		$S_{3}$ & $1010/0000$ & $0011/0011$ & $1010/0000$ & $1010/0000$ \\
		$S_{4}$ & $0101/0100$ & $0110/0101$ & $0000/0000$ & $0000/0000$ \\
		$S_{5}$ & $0100/0100$ & $0100/0100$ & $0100/0100$ & $0100/0100$ \\
		$S_{6}$ & $0100/0101$ & $0100/0101$ & $0110/0110$ & $0100/0101$ \\
		$S_{7}$ & $0111/0111$ & $1000/0000$ & $1011/1010$ & $0000/0000$ \\
		$S_{8}$ & $1000/1000$ & $1001/0000$ & $1001/0000$ & $1001/0000$ \\
		$S_{9}$ & $1001/1001$ & $0111/0000$ & $0111/0000$ & $0111/0000$ \\
		$S_{10}$ & $0000/0000$ & $0000/0000$ & $0000/0000$ & $0000/0000$ \\
		$S_{11}$ & $1011/1011$ & $0111/0000$ & $0111/0000$ & $0111/0000$ \\
		\hline
	\end{tabularx}
\end{table}

